
\SGASection{8}{Representability criterion for the functor of sub-tori
of a smooth group}
\SGARefTarget{sga3:xv:section:8}

\medskip\noindent\SGAInlineRefTarget{sga3:xv:paragraph:8-0}\textbf{8.0.}
In this section, if \(S\) is a prescheme and \(G\) an \(S\)-group
prescheme, \(\mathcal T_G\) (or simply \(\mathcal T\), if there is no
ambiguity) denotes the \(S\)-functor
\[
  (\operatorname{Sch}/S)^\circ\longrightarrow\operatorname{Ens}
\]
such that, for every \(S\)-prescheme \(S'\),
\[
  \mathcal T(S')=\{\text{sub-tori of \(G_{S'}\)}\}.
\]
The functor \(\mathcal T_C\) of central sub-tori of \(G\) is defined
similarly.
\SGARefTarget{sga3:xv:paragraph:8:0}

\SGAStatement{\SGAInlineRefTarget{sga3:xv:proposition:8-1}Proposition 8.1}{Let \(S\) be a locally noetherian
prescheme and let \(G\) be an \(S\)-group prescheme of finite type.
Then the following conditions are equivalent:
\begin{enumerate}
\item[\textup{i)}] The functor \(\mathcal T_G\) ``commutes with adic
limits of local artinian rings'': for every \(S\)-prescheme
\(S'=\operatorname{Spec}(A)\), where \(A\) is a complete noetherian
local ring for the topology defined by its maximal ideal
\(\mathfrak m\), the canonical map
\[
  \mathcal T(S')\longrightarrow
  \varprojlim_n\mathcal T(S'_n),
  \qquad S'_n=\operatorname{Spec}(A/\mathfrak m^n),
\]
is bijective.
\item[\textup{ii)}] The assertion in \textup{i)}, restricted to the
case in which \(A\) is a complete discrete valuation ring with
algebraically closed residue field.
\item[\textup{iii)}] The assertion in \textup{ii)}, restricted to the
subfunctor \(\mathcal T^{(1)}\) of \(\mathcal T\) formed by sub-tori
of \(G\) of relative dimension \(1\).
\item[\textup{i bis)}] For every \(S\)-prescheme \(S'\) as in
\textup{i)} and every \(S'\)-torus \(T\), the canonical map
\[
  \operatorname{Hom}_{S'\text{-gr}}(T,G_{S'})
  \longrightarrow
  \varprojlim_n
  \operatorname{Hom}_{S'_n\text{-gr}}(T_{S'_n},G_{S'_n})
\]
is bijective.
\item[\textup{ii bis)}] The assertion in \textup{i bis)}, restricted
to the case in which \(A\) is a complete discrete valuation ring
with algebraically closed residue field.
\item[\textup{iii bis)}] The assertion in \textup{ii bis)}, restricted
to \(T=\mathbb G_{m,S'}\).
\end{enumerate}}
\SGARefTarget{sga3:xv:proposition:8:1}

\medskip\noindent\textbf{Remark 8.2.}
There is an analogous proposition in which one restricts to central
sub-tori of \(G\) and to central homomorphisms from a torus into \(G\).
\SGARefTarget{sga3:xv:remark:8:2}

\medskip\noindent\textit{Proof.}
We shall use the following lemma.

\SGAStatement{Lemma 8.3}{Let \(S\) be a prescheme, \(G\) an
\(S\)-group prescheme, \(T\) an \(S\)-torus, and
\(u:T\to G\) an \(S\)-group morphism. Suppose, moreover, either that
\(G\) is of finite presentation over \(S\), or that \(S\) is locally
noetherian and \(G\) is locally of finite type. Then:
\begin{enumerate}
\renewcommand{\labelenumi}{\textup{\alph{enumi})}}
\item \(\ker u\) is a subgroup of multiplicative type of \(T\);
\item the quotient \(T'=T/\ker u\) is a torus;
\item the canonical monomorphism \(T'\to G\) induced by \(u\) after
passage to the quotient is an immersion.
\end{enumerate}}
\SGARefTarget{sga3:xv:lemma:8:3}

When \(G\) is separated over \(S\), this lemma follows from
Exposé \SGARefLink{sga3:ix:theorem:6:8}{IX~6.8}. In the general case, the usual reduction allows us to
suppose that \(S\) is noetherian. We first show that
\(K=\ker u\) is flat. We may suppose that \(S\) is the spectrum of a
local artinian ring (EGA \(0_{\mathrm{III}}\), 10.2.6); then \(G\) is
separated (Exposé VI\textsubscript{B}~5.2), so \(K\) is of
multiplicative type (Exposé \SGARefLink{sga3:ix:theorem:6:8}{IX~6.8}), and hence is flat over \(S\).
We next show that \(\ker u\) is closed in \(T\). This reduces us to
the case in which \(S\) is the spectrum of a discrete valuation ring
(EGA II~7.2.1). Since \(T\) is flat with connected fibres, \(u\)
factors through the identity component of the schematic closure in
\(G\) of the generic fibre of \(G\). We may thus suppose that \(G\)
is flat with connected fibres; it is then separated
(Exposé VI\textsubscript{B}~\SGARefLink{sga3:xv:theorem:5-2}{5.2}), and \(K\) is closed.
It follows from Exposé \SGARefLink{sga3:x:corollary:4-8:item:b}{X~4.8(b)} that \(K\) is a subgroup of
multiplicative type of \(T\). The quotient \(T'=T/K\) is therefore
representable and is a group of multiplicative type
(Exposé \SGARefLink{sga3:ix:proposition:2:7}{IX~2.7}) whose fibres are tori; hence \(T'\) is a torus.
Finally, the monomorphism \(T'\to G\) is an immersion by
Exposé \SGARefLink{sga3:viii:proposition:7:9}{VIII~7.9}.

\medskip\noindent\textit{Proof of \SGARefLink{sga3:xv:proposition:8-1}{Proposition 8.1}.}
\textup{i)}\(\Rightarrow\)\textup{i bis)}.
Put \(T_n=T_{S'_n}\), \(G_n=G_{S'_n}\), and let
\[
 (u_n)_{n\in\mathbb N}\in
 \varprojlim_n\operatorname{Hom}_{S'_n\text{-gr}}(T_n,G_n).
\]
For every \(n\), \(u_n(T_n)\) is a sub-torus \(T'_n\) of \(G_n\)
(\SGARefLink{sga3:xv:lemma:8:3}{Lemma 8.3}). By hypothesis, there is a unique sub-torus \(T'\) of
\(G\) lifting \(T'_n\) for every \(n\). Since \(T'\) has affine
fibres, the assertion follows from \SGARefLink{sga3:xv:corollary:4:4}{4.4}.

\textup{i bis)}\(\Rightarrow\)\textup{i)}.
Let
\((T_n)_{n\in\mathbb N}\in\varprojlim_n\mathcal T(S'_n)\).
By Exposé \SGARefLink{sga3:x:corollary:4-6}{X~4.6} there are an \(S'\)-torus \(T^0\) and an
\(S'_0\)-isomorphism
\[
  u_0:T^0_0\xrightarrow{\sim}T_0.
\]
Since \(T_n\) is smooth over \(S'_n\), \(u_0\) lifts to an
\(S'_n\)-morphism \(u_n:T^0_n\to T_n\) (Exposé \SGARefLink{sga3:ix:theorem:3:6}{IX~3.6}), and the
family \((u_n)_{n\in\mathbb N}\) may be chosen coherent. It therefore
comes from a morphism \(u:T^0\to G\). By \SGARefLink{sga3:xv:lemma:8:3}{Lemma 8.3}, the image of
\(T^0\) under \(u\) is a sub-torus of \(G\), and it lifts \(T_n\)
for every \(n\).

The implications
\[
 \textup{i)}\Rightarrow\textup{ii)}\Rightarrow\textup{iii)}
 \quad\text{and}\quad
 \textup{i bis)}\Rightarrow\textup{ii bis)}
 \Rightarrow\textup{iii bis)}
\]
are immediate. The implication
\(\textup{iii)}\Rightarrow\textup{iii bis)}\) is proved in the same
way as
\(\textup{i)}\Rightarrow\textup{i bis)}\).
It remains to prove
\[
 \textup{iii bis)}\Rightarrow\textup{ii bis)}
 \Rightarrow\textup{i bis)}.
\]

\textup{ii bis)}\(\Rightarrow\)\textup{i bis)}.
In the terminology of \SGARefLink{sga3:xv:proposition:4:3}{4.3}, assertion \textup{i bis)} holds if and
only if every element
\[
 (u_n)_{n\in\mathbb N}\in
 \varprojlim_n\operatorname{Hom}_{S'_n\text{-gr}}(T_n,G_n)
\]
is admissible. For every point \(t\) of \(S'\) distinct from its
closed point \(s\), there is an \(S\)-scheme \(S''\), the spectrum of
a complete discrete valuation ring with algebraically closed residue
field, whose generic point maps to \(t\) and whose closed point maps
to \(s\) (EGA II~7.1.9). This gives at once a valuative criterion for
a family of morphisms to be admissible, and proves the implication.

\textup{iii bis)}\(\Rightarrow\)\textup{ii bis)}.
Let \(T\) be an \(S\)-torus, where \(S\) is the spectrum of a complete
discrete valuation ring with algebraically closed residue field. The
torus \(T\) is split (Exposé \SGARefLink{sga3:x:corollary:4-6}{X~4.6}), say
\[
  T\simeq\mathbb G_{m,S}^{\,r}.
\]
Let
\[
 (u_n)_{n\in\mathbb N}\in
 \varprojlim_n
 \operatorname{Hom}_{S_n\text{-gr}}
 \bigl((\mathbb G_m)^r_{S_n},G_n\bigr).
\]
By hypothesis, the restriction of \(u_n\) to each factor of
\(\mathbb G_{m,S}^{\,r}\) comes from a group morphism
\(\mathbb G_{m,S}\to G\). Their product, followed by the composition
law \(G^r\to G\), gives a morphism
\[
  v:\mathbb G_{m,S}^{\,r}\longrightarrow G.
\]
Clearly \(u_n=v_n\) for every \(n\). It remains to prove that \(v\)
is a group morphism. Equivalently, two evident morphisms
\[
 f,g:X=\mathbb G_{m,S}^{\,r}\times_S\mathbb G_{m,S}^{\,r}
 \longrightarrow G
\]
must coincide. Let \(Z\) be their coincidence subscheme. Since
\(\mathbb G_{m,S}^{\,r}\) has connected fibres and its image under
\(v\) contains the unit section, the argument of \SGARefLink{sga3:xv:lemma:8:3}{Lemma 8.3} shows
that \(v\) factors through the identity component of \(G\). We may
therefore suppose that \(G\) is separated
(Exposé VI\textsubscript{B}~\SGARefLink{sga3:xv:theorem:5-2}{5.2}), so \(Z\) is closed. On the other
hand, \(v_n\) is a group morphism, hence \(f_n=g_n\) for every \(n\).
Thus \(Z\) contains a neighbourhood of the closed fibre of \(X\)
(EGA I~10.9.4), which is schematically dense in \(X\), since \(X\)
has smooth irreducible fibres (Exposé IX~4.6). Consequently
\(Z=X\), and \(f=g\).

\SGAStatement{\SGAInlineRefTarget{sga3:xv:definition:8-4}Definition 8.4}{Let \(S\) be a locally noetherian
prescheme and \(G\) an \(S\)-group prescheme of finite type. Let
\(\mathcal S\) be the set of \(S\)-schemes \(S'\) that are spectra
of complete discrete valuation rings with algebraically closed
residue field. Write \(s\) and \(t\) for the closed and generic
points of \(S'\), and require that
\((G_t)^0_{\mathrm{red}}\) be smooth and admit a Chevalley
decomposition: an extension of an abelian variety by a smooth
connected linear algebraic subgroup \(L_t\). This decomposition is
then unique. For \(S'\in\mathcal S\), let \(G'\), respectively \(L\),
denote the schematic closure in \(G_{S'}\) of \(G_t\), respectively
of \(L_t\).

We say that ``the abelian part of \(G\) does not degenerate into a
toric part'', or briefly that \(G\) has property \(\mathrm{AT}\), if,
for every \(S'\in\mathcal S\), \(L_s\) has the same reductive rank
as \(G'_s\). Intuitively, suppose that the quotient \(A=G'/L\) is
representable. Then \(A\) is a flat group prescheme and
\((A_t)^0_{\mathrm{red}}\) is an abelian variety; property
\(\mathrm{AT}\) says that \(A_s\) has reductive rank zero, so that
\((A_s)^0_{\mathrm{red}}\) is an extension of an abelian variety by
a unipotent group.

Similarly, if \(G\) has connected fibres, we say that \(G\) has
property \(\mathrm{ATC}\) if, for every \(S'\in\mathcal S\), the
schematic closure \(Z\) of the centre \(Z_t\) of \(G_t\) has
property \(\mathrm{AT}\).}
\SGARefTarget{sga3:xv:definition:8:4}

These technical definitions are justified by the following
proposition.

\SGAStatement{Proposition 8.5}{Let \(S\) be a locally noetherian
prescheme and let \(G\) be an \(S\)-group prescheme of finite type.
\begin{enumerate}
\renewcommand{\labelenumi}{\textup{\alph{enumi})}}
\item\SGAInlineRefTarget{sga3:xv:proposition:8-5:item:a} The functor \(\mathcal T\) of sub-tori of \(G\) commutes with
adic limits of local artinian rings (cf. \SGARefLink{sga3:xv:proposition:8-1}{Proposition 8.1}) if and
only if \(G\) has property \(\mathrm{AT}\).
\item\SGAInlineRefTarget{sga3:xv:proposition:8-5:item:b} If \(G\) has connected fibres, the functor \(\mathcal T_C\) of
central sub-tori of \(G\) commutes with adic limits of local artinian
rings if and only if \(G\) has property \(\mathrm{ATC}\).
\end{enumerate}}
\SGARefTarget{sga3:xv:proposition:8:5}

We need the following technical lemma.

\SGAStatement{Lemma 8.6}{Let \(S\) be the spectrum of a discrete
valuation ring, with closed point \(s\), let \(G\) be an
\(S\)-group prescheme flat and of finite type, and let \(T_s\) be a
sub-torus of \(G_s\). There are an \(S\)-scheme \(S'\), the spectrum
of a discrete valuation ring faithfully flat over \(S\), with closed
point \(s'\), and a subgroup scheme \(C\) of \(G_{S'}\), flat over
\(S'\), commutative, and with connected fibres, such that
\(C_{s'}\) contains \(T_s\times_s s'\).}
\SGARefTarget{sga3:xv:lemma:8:6}

After replacing \(S\) by the spectrum \(S'\) of a discrete valuation
ring faithfully flat over it, we may suppose that
\(T_s=\mathbb G_{m,s}^{\,r}\) and that the transcendence degree of
\(\kappa(s)\) over its prime field is at least \(r\)
(EGA \(0_{\mathrm{III}}\), 10.3.1, and EGA II~7.1.9). There is then
an element \(x\in T_s(s)\) such that every algebraic subgroup of
\(G_s\) containing \(x\) contains \(T_s\)
(cf. Exposé XIII, proof of \SGARefLink{sga3:xiii:theorem:2:1}{2.1},
\(\textup{ii)}\Rightarrow\textup{vii)}\)).
Since \(G\) is flat over \(S\), a quasi-section passes through \(x\)
(EGA IV~14.5.8). After another such replacement of \(S\), we may
therefore suppose that \(x\) lifts to a section
\(\widetilde x\) of \(G\) over \(S\). Let \(C_t\) be the commutative
algebraic subgroup of \(G_t\) generated by \(\widetilde x_t\)
(Exposé VI\textsubscript{B}~7), and let \(C\) be its schematic
closure in \(G\). Clearly \(C_s\) contains \(x\), hence contains
\(T_s\). The identity component of \(C\) is therefore a flat
commutative group scheme with connected fibres having the required
property.

\medskip\noindent\textit{Proof of \SGARefLink{sga3:xv:proposition:8-5:item:a}{Proposition 8.5(a)}.}
Suppose first that \(\mathcal T\) commutes with adic limits of
artinian rings. Let \(S'\in\mathcal S\), and let \(T_s\) be a maximal
torus of \(G'_s\). We must show that \(T_s\subset L_s\). The
formation of \(L\) and \(G'\) plainly commutes with faithfully flat
extensions \(S''\to S'\) of discrete valuation rings. By \SGARefLink{sga3:xv:lemma:8:6}{Lemma 8.6},
we may therefore replace \(S'\) so that there is a flat commutative
subgroup scheme \(C\) of \(G'_{S'}\) for which \(C_s\) contains
\(T_s\). The torus \(T_s\) is central in \(C_s\), and hence lifts
infinitesimally to a central sub-torus (\SGARefLink{sga3:xv:corollary:2:3}{Corollary 2.3}). By the
hypothesis on \(G\), \(T_s\) lifts to a sub-torus \(T\) of \(G\).
Clearly \(T_t\) lies in the linear part \(L_t\) of \(G_t\), and
therefore \(T\subset L\).

Conversely, suppose that \(G\) has property \(\mathrm{AT}\). We prove
condition \SGARefLink{sga3:xv:proposition:8-1}{8.1} \textup{iii bis)}. Let \(S=\operatorname{Spec}(A)\),
where \(A\) is a complete discrete valuation ring with algebraically
closed residue field, let \(\mathfrak m\) be its maximal ideal, put
\(S_n=\operatorname{Spec}(A/\mathfrak m^n)\), and let
\[
 u_n:(\mathbb G_{m,S})_n\longrightarrow
 G_n=G\times_S S_n,\qquad n\in\mathbb N,
\]
be a coherent system of group morphisms. Let \(q\) be a prime
invertible on \(S\). As \(\ell\) ranges over the powers of \(q\),
there is a unique \(S\)-group morphism
\[
 {}_\ell u:{}_\ell\mathbb G_{m,S}\longrightarrow G
\]
lifting the restriction of \(u_n\) to
\({}_\ell(\mathbb G_{m,S})_n\) for every \(n\)
(\SGARefLink{sga3:xv:proposition:1-6:item:a}{Proposition 1.6(a)}). Consequently, if an \(S\)-group morphism
\(u:\mathbb G_{m,S}\to G\) lifting every \(u_n\) exists, its
restriction to \({}_\ell\mathbb G_{m,S}\) is uniquely determined.
The density theorem (Exposé \SGARefLink{sga3:ix:corollary:4:8:item:a}{IX~4.8(a)}) proves both the uniqueness of
\(u\) and that, for its existence, we may make a faithfully flat
extension of the base.

Now \((G_t)^0_{\mathrm{red}}\) admits a Chevalley decomposition,
already defined over a finite extension \(L\) of the fraction field
\(K\) of \(A\). After replacing \(S\) by its normalization in \(L\),
we may suppose that \(S\in\mathcal S\). The morphism
\({}_\ell u_t\) factors through \((G_t)_{\mathrm{red}}\), and
\({}_\ell u\) therefore factors through the schematic closure
\(G''\) of \((G_t)_{\mathrm{red}}\) in \(G'\). The density theorem
then shows that every \(u_n\) factors through \(G''_n\).
Since \(G\) has property \(\mathrm{AT}\), every sub-torus of \(G'_s\),
and hence of \(G''_s\), is contained in \(L_s\); thus \(u_s\)
factors through \(L_s\). In fact every \(u_n\) factors through
\(L_n\). Indeed, \(L_t\) is invariant in \(G''_t\), so \(L\) is
invariant in \(G''\); moreover \(L\) is flat over \(S\), and the
quotient
\[
  H_n=G''_n/L_n
\]
is representable for every \(n\) (Exposé VI\textsubscript{A},
\S4). The image of \((\mathbb G_{m,S})_n\) in \(H_n\) is a
sub-torus (Exposé \SGARefLink{sga3:ix:theorem:6:8}{IX~6.8}) whose closed fibre is zero. It is therefore
zero, and \(u_n\) factors through \(L_n\). Since \(L_t\) is affine,
\SGARefLink{sga3:xv:proposition:4:3}{Proposition 4.3} now shows that the family \((u_n)\) is admissible.

The proof of \SGARefLink{sga3:xv:proposition:8-5:item:b}{Proposition 8.5(b)} is entirely analogous, in view of
\SGARefLink{sga3:xv:remark:8:2}{Remark 8.2} and Exposé \SGARefLink{sga3:ix:theorem:5:6:item:a}{IX~5.6(a)}.

\SGAStatement{Proposition 8.7}{Let \(S\) be a locally noetherian
prescheme and let \(G\) be an \(S\)-group prescheme of finite type.
\begin{enumerate}
\renewcommand{\labelenumi}{\textup{\roman{enumi})}}
\item If \(G\) has connected fibres and has property
\(\mathrm{AT}\), then it has property \(\mathrm{ATC}\).
\item If \(G\) has property \(\mathrm{AT}\), every subgroup
prescheme of \(G\) has property \(\mathrm{AT}\).
\item If \(G\) is flat over \(S\) and the abelian rank of its fibres
is a locally constant function on \(S\), then \(G\) has property
\(\mathrm{AT}\).
\end{enumerate}}
\SGARefTarget{sga3:xv:proposition:8:7}

Assertion \textup{(i)} follows at once from \SGARefLink{sga3:xv:proposition:8:5}{Proposition 8.5} and
Exposé \SGARefLink{sga3:ix:theorem:5:6:item:a}{IX~5.6(a)}.

For \textup{(ii)}, let \(H\) be a subgroup prescheme of \(G\). To
prove that \(H\) has property \(\mathrm{AT}\), we may suppose that
\(S\) is the spectrum of a complete discrete valuation ring. We
must show that every coherent family of \(S_n\)-group morphisms
\[
  u_n:(\mathbb G_{m,S})_n\longrightarrow H_n
\]
is admissible (\SGARefLink{sga3:xv:proposition:8:5}{Proposition 8.5} and \SGARefLink{sga3:xv:proposition:8-1}{8.1} \textup{iii bis)}).
By hypothesis, \(G\) has property \(\mathrm{AT}\), so there is an
\(S\)-group morphism
\[
  u:\mathbb G_{m,S}\longrightarrow G
\]
lifting every \(u_n\). As in the proof of \SGARefLink{sga3:xv:proposition:8:5}{Proposition 8.5}, the
restriction of \(u\) to \({}_\ell\mathbb G_{m,S}\), with \(\ell\)
ranging over powers of a prime \(q\) invertible on \(S\), factors
through \(H\). By density, \(u_t\) factors through \(H_t\) on the
generic fibre. Since \(\mathbb G_{m,S}\) is reduced, \(u\) itself
factors through \(H\).

For \textup{(iii)}, let \(S'\in\mathcal S\), with the notation of
\SGARefLink{sga3:xv:definition:8-4}{Definition 8.4}. The group \(L\) is flat over \(S\), and its generic
fibre \(L_t\) is affine. Under these conditions \(L_s\) is
necessarily affine (Exposé XVII, Appendix III~2). Since \(G\) and
\(L\) are flat, \(G=G'\), and the dimensions of the fibres of \(G\)
and \(L\) are constant on \(S\) (Exposé VI\textsubscript{B}~4).
Consequently
\[
 \operatorname{rk}_{\mathrm{ab}}G_s
 \leq \dim G_s-\dim L_s
 =\dim G_t-\dim L_t
 =\operatorname{rk}_{\mathrm{ab}}G_t.
\]
By hypothesis these inequalities are equalities. Thus
\((G_s/L_s)^0_{\mathrm{red}}\) is an abelian variety and has
reductive rank zero. It follows that \(L_s\) has the same reductive
rank as \(G_s\).

We can now prove the principal theorems of this section.

\SGAStatement{\SGAInlineRefTarget{sga3:xv:theorem:8-8}Theorem 8.8}{Let \(G\) be a group prescheme of finite
type over a locally noetherian prescheme \(S\). Suppose that \(G\)
is flat over \(S\) and has connected fibres.
\begin{enumerate}
\renewcommand{\labelenumi}{\textup{\alph{enumi})}}
\item\SGAInlineRefTarget{sga3:xv:theorem:8-8:item:a} The functor \(\mathcal T_C\) of central sub-tori of \(G\) is
representable if and only if \(G\) has property \(\mathrm{ATC}\).
In that case, \(\mathcal T_C\) is represented by an \(S\)-prescheme
étale and separated over \(S\).
\item\SGAInlineRefTarget{sga3:xv:theorem:8-8:item:b} Under the hypotheses of \textup{(a)}, for every \(S\)-torus
\(T\), the functor
\(\operatorname{Hom}^{\mathrm{cent}}_{S\text{-gr}}(T,G)\) of central
homomorphisms from \(T\) into \(G\) is represented by an
\(S\)-prescheme étale and separated over \(S\).
\end{enumerate}}
\SGARefTarget{sga3:xv:theorem:8:8}

\SGAStatement{Theorem 8.9}{Let \(S\) be a locally noetherian
prescheme and let \(G\) be an \(S\)-group prescheme of finite type,
smooth over \(S\).
\begin{enumerate}
\renewcommand{\labelenumi}{\textup{\alph{enumi})}}
\item\SGAInlineRefTarget{sga3:xv:theorem:8-9:item:a} The functor \(\mathcal T\) of sub-tori of \(G\) is
representable if and only if \(G\) has property \(\mathrm{AT}\).
In that case, \(\mathcal T\) is represented by an \(S\)-prescheme
smooth and separated over \(S\). More precisely, its structural
morphism admits a canonical factorization
\[
  \mathcal T\xrightarrow{u}Y\xrightarrow{v}S,
\]
where \(v\) is smooth and quasi-projective (hence of finite type)
and \(u\) is étale and separated.
\item\SGAInlineRefTarget{sga3:xv:theorem:8-9:item:b} Under the hypotheses of \textup{(a)}, for every \(S\)-torus
\(T\), the functor \(\operatorname{Hom}_{S\text{-gr}}(T,G)\) of
homomorphisms from \(T\) into \(G\) is represented by an
\(S\)-prescheme smooth and separated over \(S\).
\end{enumerate}}
\SGARefTarget{sga3:xv:theorem:8:9}
